\chapter{Identifying the perimeter measure}\label{sec:perimeter-measure-identification}
The graph cover and the pointwise density limit still have to be connected to the perimeter measure. We separate the needed measure comparisons and state the remaining differentiation input explicitly. This completes the structure proof from \cite{Sim18}, Chapter 3, Theorem 4.4, printed pp.~97--100, relative to the Radon differentiation theorem proved from the cited Besicovitch covering input and the rectifiable-set density theorem proved through $C^1$ pieces.

\begin{corollary}[The perimeter measure is concentrated on the reduced boundary]\label{cor:perimeter-concentrated-reduced-boundary}
For a locally finite-perimeter set $E$ in $U$,
\[
 |D\chi_E|(U\setminus\partial^*E)=0.
\]
This is the concentration assertion at the beginning of the proof of \cite{Sim18}, Chapter 3, Theorem 4.4, p.~97.
\end{corollary}
\begin{proof}
Put $\mu=|D\chi_E|$ and write $D\chi_E=-\nu\mu$ with $|\nu|=1$ almost everywhere. The bounded polar field belongs locally to $L^1(\mu)$. Theorem \ref{thm:radon-function-differentiation}, applied to its finitely many components, gives
\[
 \frac{D\chi_E(B_r(x))}{\mu(B_r(x))}\longrightarrow-\nu(x)
 \quad\text{for }\mu\text{-almost every }x.
\]
The complement of the support of $\mu$ is a countable union of open balls with zero measure, so almost every such point also lies in that support. Since the limit has length one, these points belong to the reduced boundary by definition. Its Borel measurability follows from Lemma \ref{lem:reduced-boundary-borel}, so the displayed measure statement is well defined.
\end{proof}

\begin{proposition}[A lower density bound controls Hausdorff measure]\label{prop:lower-density-hausdorff-comparison}
Let $\mu$ be a positive Radon measure on $U\subset\mathbb R^n$, let $m>0$, and let $S\subset U$ be Borel. Suppose that for each $x\in S$ and every sufficiently small $r>0$,
\[
 \mu(B_r(x))\geq c r^m
\]
with a fixed $c>0$; the permitted scale may depend on $x$. Then for every Borel $A\subset S$,
\[
 \mathcal H^m(A)\leq \frac{5^m\omega_m}{c}\mu(A).
\]
Here $m$ is an integer between one and $n$ so that $\omega_m$ has the normalization of Definition \ref{def:hausdorff-measure}. The bound is not sharp.
\end{proposition}
\begin{proof}
First suppose $A$ is bounded and lies in a compact subset of $U$. Take an open $O\subset U$ containing $A$ and a positive diameter scale $\delta$. For each $x\in A$ use all sufficiently small balls centered at $x$ that lie in $O$, have radius less than $\delta/10$, and satisfy the stated mass lower bound. This family covers $A$ and has bounded centers and radii. Lemma \ref{lem:bounded-fivefold-covering} provides disjoint selected balls $B_{r_i}(x_i)$ whose fivefold dilations cover $A$. Their diameters are less than $\delta$, so the normalized Hausdorff content satisfies
\[
 \mathcal H^m_\delta(A)
 \leq\sum_i\omega_m(5r_i)^m
 \leq\frac{5^m\omega_m}{c}\sum_i\mu(B_{r_i}(x_i))
 \leq\frac{5^m\omega_m}{c}\mu(O).
\]
Let $\delta\downarrow0$ and then use outer regularity of $\mu$ to take the infimum over $O$. It is enough to use open sets inside a fixed relatively compact neighborhood of $A$, so their measures are finite in this step. Finally exhaust $U$ by compact sets and apply the result to their intersections with $A$. These sets increase to $A$, and both Borel measures are continuous from below, proving the general assertion. The constants come from fivefold dilation and the normalized diameter definition, not a sharp density comparison theorem.
\end{proof}

\begin{proposition}[An upper density bound controls the restricted Radon measure]\label{prop:upper-density-radon-comparison}
Let $\mu$ be a positive Radon measure on $U$, and let $S\subset U$ be Borel. Suppose that for every $x\in S$ and all sufficiently small $r>0$,
\[
 \mu(B_r(x))\leq C\omega_m r^m,
\]
where $1\leq m\leq n$ and $C$ is fixed, but the permitted scale may depend on $x$. Then for every Borel $A\subset S$,
\[
 \mu(A)\leq C4^m\mathcal H^m(A).
\]
This is the diameter-cover comparison used after equation (13) in \cite{Sim18}, p.~100, with a convenient enlarged radius.
\end{proposition}
\begin{proof}
Let $S_j$ consist of points of $S$ for which $\operatorname{dist}(x,\mathbb R^n\setminus U)>1/j$ and the stated upper bound holds for all $0<r<1/j$. These are Borel sets: it suffices to test positive rational radii, because continuity from below of open balls passes the bound to every real radius. They increase with $j$ and cover $S$.

Fix $j$ and cover $A\cap S_j$ by sets $Q_i$ of diameters $d_i<\delta<1/(2j)$, with total normalized cost within $\varepsilon$ of $\mathcal H^m_\delta(A\cap S_j)$. Discard sets missing $A\cap S_j$, and choose $x_i\in Q_i\cap A\cap S_j$ for the others. When $d_i>0$, $Q_i\subset B_{2d_i}(x_i)$, and
\[
 \mu(B_{2d_i}(x_i))\leq C\omega_m(2d_i)^m
 =C4^m\omega_m(d_i/2)^m.
\]
When $d_i=0$, $Q_i$ contains at most the single point $x_i$, and $\mu(\{x_i\})=0$ by the assumed upper bound and $m>0$. Subadditivity gives
\[
 \mu(A\cap S_j)\leq C4^m
             \bigl(\mathcal H^m_\delta(A\cap S_j)+\varepsilon\bigr)
 \leq C4^m\bigl(\mathcal H^m(A)+\varepsilon\bigr).
\]
If the Hausdorff measure on the right is infinite the result is immediate; otherwise such finite-cost covers exist. Let $\varepsilon\downarrow0$ and then $j\to\infty$, using continuity from below of $\mu$. This proves the assertion without assuming a uniform radius on all of $A$.
\end{proof}

\begin{corollary}[Perimeter and reduced-boundary Hausdorff measure have the same null sets]\label{cor:perimeter-hausdorff-equivalence}
For $E$ of locally finite perimeter in $U$, put $m=n-1\geq1$, $S=\partial^*E\cap U$, and $\sigma=\mathcal H^m\llcorner S$. Then $\sigma$ is a locally finite Radon measure and
\[
 \sigma\leq2\cdot5^m|D\chi_E|,
 \qquad |D\chi_E|\leq2\cdot4^m\sigma.
\]
These coarse comparisons yield mutual absolute continuity. The exact equality requires the next step.
\end{corollary}
\begin{proof}
The density limit in Corollary \ref{cor:reduced-boundary-exact-densities} gives, at each point of $S$ and at all sufficiently small radii, the lower bound $(\omega_m/2)r^m$ and the upper bound $2\omega_m r^m$. Apply Proposition \ref{prop:lower-density-hausdorff-comparison} with $c=\omega_m/2$, and Proposition \ref{prop:upper-density-radon-comparison} with $C=2$. The latter initially bounds $|D\chi_E|\llcorner S$; Corollary \ref{cor:perimeter-concentrated-reduced-boundary} identifies this with the whole perimeter measure inside $U$. The Borel set $S$ is measurable for Hausdorff measure, and the first bound shows local finiteness of its restriction. A locally finite Borel measure on a Euclidean open set is Radon. Thus the comparisons are inequalities of Radon measures.
\end{proof}

\begin{theorem}[Unit density identifies the perimeter measure exactly]\label{thm:perimeter-hausdorff-identification}
For $E$ of locally finite perimeter in $U\subset\mathbb R^n$, $n\geq2$,
\[
 |D\chi_E|=\mathcal H^{n-1}\llcorner(\partial^*E\cap U).
\]
The proof uses the arbitrary-Radon differentiation theorem proved from Besicovitch selection and the locally proved density theorem for rectifiable sets, in addition to the local arguments developed above. This is the final measure-identification step of \cite{Sim18}, Theorem 4.4, p.~100.
\end{theorem}
\begin{proof}
Put $\mu=|D\chi_E|$, $S=\partial^*E\cap U$ and $\sigma=\mathcal H^{n-1}\llcorner S$. Corollary \ref{cor:perimeter-hausdorff-equivalence} gives locally finite mutually absolutely continuous measures. The abstract Radon--Nikodym theorem gives $\mu=f\sigma$ for a nonnegative locally $\sigma$-integrable function $f$; one may apply it on a countable finite-measure exhaustion and use uniqueness on overlaps. This abstract representation is recalled in \cite{Sim18}, Chapter 1, equation (4.17), p.~27.

By Theorem \ref{thm:radon-function-differentiation},
\[
 f(x)=\lim_{r\downarrow0}\frac{\mu(B_r(x))}{\sigma(B_r(x))}
 \quad\text{for }\sigma\text{-almost every }x.
\]
Theorem \ref{thm:reduced-boundary-graph-cover} makes $S$ countably rectifiable, and the preceding corollary supplies the required local Hausdorff finiteness. Corollary \ref{cor:rectifiable-local-unit-density} therefore gives $\sigma(B_r(x))/(\omega_{n-1}r^{n-1})\to1$ at $\sigma$-almost every point. At every point of $S$, Corollary \ref{cor:reduced-boundary-exact-densities} gives the same unit limit for $\mu$. Dividing the two limits yields $f=1$ almost everywhere for $\sigma$, proving the measure identity. The density theorem is used only after its rectifiability and local-finiteness hypotheses have been established, so this step does not assume the perimeter representation it is proving.
\end{proof}
